\paragraph{Generic vs.\ type-specific prompts.}
We additionally tested a generic disambiguation prompt (``Before acting, check whether the policy clause could be interpreted in more than one way; if so, ask the customer for clarification'') against the type-specific prompts reported above, using GPT-5.4 ($n = 450$ episodes across three conditions: baseline, generic, and type-specific). Table~\ref{tab:generic_vs_specific} reports the results. The generic prompt produced no significant reduction in any type (overall $\Delta = +1.3$\,pp, all per-type McNemar $p > 0.05$), consistent with the type-dependent effectiveness pattern: authorization scope showed the largest (non-significant) reduction under both generic ($-$6.0\,pp) and type-specific ($-$12.0\,pp) prompts, while incompleteness showed a non-significant \emph{increase} under the type-specific prompt ($+$14.0\,pp, $p = 0.118$). This single-model pilot has limited statistical power ($n = 50$ per cell); multi-model validation is left to future work. The pattern nevertheless reinforces the main finding that prompt-level intervention is insufficient for specification-layer ambiguity, particularly incompleteness, where the missing information cannot be recovered from the prompt alone.
